\section{Introduction}
\label{sec:introduction}

\IEEEPARstart{A}{pproximate}  computing trades bounded degradation in output quality for improvements in hardware area, power, and delay, providing an effective means of improving the efficiency of domain-specific accelerators~\cite{jiang2020approxarith,que2023approxsurvey,xu2016approxsurvey}. 
Approximate arithmetic libraries, such as EvoApprox8b~\cite{mrazek2017evoapprox}, provide diverse adders, subtractors, and multipliers with different error characteristics and hardware costs, enabling designers to construct approximate accelerators by selectively replacing exact arithmetic units.
However, when an accelerator contains many configurable arithmetic sites and each site admits multiple legal implementations, the number of system configurations grows combinatorially.
System-level power, performance, and area (PPA) depend on logic synthesis, structural interactions, and timing paths, so component costs alone do not determine the accelerator implementation.
Application quality is likewise determined by how local approximation errors propagate through the complete computation.
Consequently, exhaustive synthesis and application-level evaluation quickly become impractical, making efficient design space exploration (DSE) a central challenge in automated approximate-accelerator design~\cite{saeedi2024dse,silva2024mapping}.

Surrogate-assisted DSE has substantially reduced the cost of exploring such large spaces.
AutoAx~\cite{mrazek2019autoax} established an automated exploration methodology based on libraries of approximate components, while Xel-FPGAs~\cite{prabakaran2023xelfpgas} extended end-to-end automated exploration to FPGA-based approximate accelerators.
More recently, ApproxPilot~\cite{zhang2025approxpilot} introduced graph-neural-network-based modeling for accelerator approximation, and ApproxGNN~\cite{vlcek2025approxgnn} explored pretrained graph representations for parameter prediction in approximate-computing DSE.
These advances demonstrate that learned surrogates can replace a large fraction of expensive hardware evaluations and enable exploration of design spaces that would otherwise be difficult to search directly.

The central challenge is to align surrogate learning with the regions that matter to optimization.
Multi-objective DSE does not query the design space uniformly.
As optimization progresses, the search increasingly concentrates on configurations that appear competitive with respect to the current Pareto front.
Therefore, prediction quality in these Pareto-relevant regions is more important to the final optimization outcome than uniformly high accuracy over the entire design space.
In representative surrogate-based approximate-accelerator flows, the surrogate is trained or adapted before the main search and is subsequently used to evaluate a much larger set of candidate configurations~\cite{mrazek2019autoax,zhang2025approxpilot,vlcek2025approxgnn}.
Evaluating promising configurations identified during DSE provides system-level supervision directly relevant to subsequent search.
Feedback-driven evaluation has already proved useful for other expensive hardware optimization problems, including Bayesian and active-learning-based DSE~\cite{reagen2017bayesopt,bai2021boomexplorer,huang2024arsflow,huang2025arsflow2}.
For approximate accelerators, the key question is how to couple such feedback with simultaneous PPA and application-quality modeling over a large, heterogeneous, multi-site approximation space.

Delay modeling poses a second challenge because its structural dependence differs from that of area and power.
While area and power reflect aggregate effects across the synthesized design, system delay is dominated by dependencies along competing combinational paths.
ApproxPilot explicitly incorporates critical-path information into its delay modeling~\cite{zhang2025approxpilot}.
More broadly, graph-based delay prediction has demonstrated the importance of structural information in both FPGA and ASIC-oriented modeling~\cite{ustun2020delay,de2023asicdelay}.
Motivated by this property, we directly propagate and aggregate latent node representations along RTL timing domains and directed dependencies.
The resulting representation captures path-dependent delay behavior at the configuration level without requiring explicit critical-path node labels.

We present \textbf{ApproxPilot-SE}, an extension of ApproxPilot that couples surrogate learning and DSE through iterative ground-truth feedback.
Separate PPA and quality-of-result (QoR) surrogates guide the search over legal configurations of the target accelerator.
Non-dominated sorting and crowding distance identify Pareto-relevant candidates for logic synthesis and application-level evaluation.
The measured labels update both surrogates, directing subsequent search toward regions informed by the new measurements.
After the iterative search terminates, the final candidate set is evaluated using logic synthesis and application-level quality measurement to construct the verified Pareto front.

The main contributions of this work are summarized as follows:

\begin{itemize}

    \item \textbf{A co-evolving PPA--QoR DSE framework for multi-site approximate-accelerator optimization.}
    ApproxPilot-SE couples multi-objective search with iterative system-level feedback: Pareto-relevant configurations identified during DSE are evaluated for real PPA and QoR, incorporated into the training set, and used to update both surrogates before the next search round.
    Under the same 900-label target-system training budget, ApproxPilot-SE improves ground-truth hypervolume (HV) gain over an aligned one-shot workflow by \textbf{50.84\%} at 99\% SSIM retention on DCT76.

    \item \textbf{A timing-domain path aggregation mechanism without explicit critical-path node supervision.}
    ApproxPilot-SE propagates latent node representations along directed RTL dependencies within timing domains and aggregates competing paths into a structure-aware representation for system-level delay prediction.
    Across three paired training seeds, the proposed readout reduces Delay MAE by \textbf{7.45\%} relative to global max pooling, while improving Spearman correlation from \textbf{0.7965} to \textbf{0.8217} and pairwise ranking accuracy from \textbf{80.37\%} to \textbf{81.67\%}.

    \item \textbf{An end-to-end, ground-truth-verified optimization flow for approximate accelerators.}
    ApproxPilot-SE integrates legal configuration generation, PPA/QoR modeling, co-evolving search, and final synthesis- and application-level verification into a unified system-level workflow.
    On DCT76, it achieves higher ground-truth PPA HV gain than ApproxPilot and ApproxGNN under all five SSIM-retention constraints from 80\% to 99\%.
    A representative high-quality design retains \textbf{99.68\%} of the Exact SSIM while reducing Area, Power, and Delay by \textbf{19.73\%}, \textbf{20.50\%}, and \textbf{12.71\%}, respectively.

\end{itemize}

The remainder of this paper is organized as follows.
Section~\ref{sec:related} reviews automated approximate-accelerator design, learning-based hardware modeling, and sequential DSE with expensive evaluations.
Section~\ref{sec:method} presents the surrogate models, timing-domain path aggregation, and the co-evolution of surrogate learning and DSE in ApproxPilot-SE.
Section~\ref{sec:experiments} evaluates ApproxPilot-SE in terms of ground-truth Pareto quality, surrogate ranking quality, the effect of co-evolving optimization, delay modeling, and representative hardware designs.
Section~\ref{sec:conclusion} concludes the paper.
